\documentclass[11pt]{article}

\usepackage[margin=1in]{geometry}
\usepackage{amsmath,amssymb}
\usepackage{booktabs}
\usepackage{hyperref}
\usepackage{xcolor}
\usepackage{enumitem}
\usepackage{listings}

\hypersetup{
  colorlinks=true,
  linkcolor=blue!60!black,
  urlcolor=blue!60!black,
  citecolor=blue!60!black
}

\newcommand{\code}[1]{\texttt{\detokenize{#1}}}
\newcommand{\R}{\mathbb{R}}

\lstset{
  basicstyle=\ttfamily\small,
  columns=fullflexible,
  keepspaces=true,
  frame=single,
  breaklines=true,
  backgroundcolor=\color{black!3}
}

\title{Level-Scheduled Upwind Block Gauss-Seidel Preconditioning}
\author{\code{hdgfem.linalg.upwind_block_gs}}
\date{\today}

\begin{document}
\maketitle

\begin{abstract}
This note documents the first implementation of a flow-aware block
Gauss-Seidel preconditioner for advection-reaction HDG trace systems.  The
method assumes the global trace matrix has already been ordered by the upwind
SCC graph ordering.  It builds a reusable SciPy \code{LinearOperator} from the
assembled CSR matrix, using dense edge-block diagonal inverses and strictly
upstream cross-level block couplings.  Couplings within the same topological
level and couplings to downstream levels are dropped, which makes each level
parallelizable with Numba.
\end{abstract}

\tableofcontents

\section{Context}

For a boundary-eliminated HDG advection-reaction trace system, the global solve
has the form
\[
  A \lambda = b,
\]
where each active mesh edge contributes
\[
  n_{\partial} = p + 1
\]
trace degrees of freedom.  After upwind SCC ordering, the unknown vector is
viewed as a sequence of edge blocks,
\[
  \lambda =
  \begin{bmatrix}
    \lambda_0 & \lambda_1 & \cdots & \lambda_{N_b-1}
  \end{bmatrix}^T,
  \qquad
  \lambda_i \in \R^{n_{\partial}}.
\]
The ordering module also computes topological level widths
\[
  w_0,\ldots,w_{L-1},
\]
where blocks in the same level are intended to be independent with respect to
the directed upwind dependency graph.

\section{Preconditioner Approximation}

Let \(A_{ij}\in\R^{n_{\partial}\times n_{\partial}}\) denote the dense block of
the trace matrix coupling block \(j\) into block row \(i\).  The version-1
preconditioner forms
\[
  M = D + L_{\rm up},
\]
where
\[
  D_{ii} = A_{ii},
\]
and \(L_{\rm up}\) contains only block couplings satisfying
\[
  \ell(j) < \ell(i),
\]
with \(\ell(i)\) the topological level of block \(i\).  All couplings with
\(\ell(j)=\ell(i)\) and \(\ell(j)>\ell(i)\) are dropped.

The default preconditioner apply computes
\[
  z = M^{-1} r
\]
by a forward level sweep.  For each level \(\ell=0,\ldots,L-1\), every block in
that level is processed in parallel:
\[
  z_i =
  A_{ii}^{-1}
  \left(
    r_i -
    \sum_{\substack{j\\ \ell(j)<\ell(i)}} A_{ij}z_j
  \right).
\]
Because all retained dependencies point from earlier levels into later levels,
the loop over levels is sequential, while the loop over blocks inside a level
is parallel.

An experimental \code{forward_backward} sweep also stores downstream
cross-level couplings.  It applies a forward sweep, multiplies by the dense
block diagonal, and then applies a backward sweep.  In matrix form, this is the
block analogue of
\[
  (D+U_{\rm down})^{-1}D(D+L_{\rm up})^{-1}.
\]
This uses both cross-level directions, but it still drops same-level couplings.
It is therefore mainly a diagnostic for whether downstream couplings are the
missing ingredient.

\section{Version-1 Builder}

The first implementation deliberately builds the preconditioner after CSR
assembly.  This gives a simple correctness path:
\begin{enumerate}[nosep]
  \item assemble the trace matrix normally;
  \item sum duplicate COO entries by converting to CSR;
  \item extract dense diagonal edge blocks;
  \item extract retained upstream block couplings into a block-CSR structure;
  \item invert every dense diagonal block;
  \item expose the Numba level sweep as a SciPy \code{LinearOperator}.
\end{enumerate}

This route avoids changing the fused trace-assembly kernel while we evaluate
the numerical quality of the preconditioner.

\section{Parallel Structure}

The apply kernel has the following structure:

\begin{lstlisting}
for level in range(num_levels):
    for block in prange(level_offsets[level], level_offsets[level + 1]):
        residual = r[block]
        for retained upstream neighbor j:
            residual -= A[block, j] @ z[j]
        z[block] = inverse_diagonal_block[block] @ residual
\end{lstlisting}

The computation is race-free because every retained neighbor \(j\) lies in an
earlier level, so \(z_j\) has already been computed before the current level
starts.  The diagnostics from the trigonometric-style \code{test2} advection
case with \(p=6\) and \(\code{lc}=0.01\) showed 439 levels and maximum width
599, which is a useful amount of level parallelism.

In practice, launching a threaded region once per level has overhead.  The
implementation therefore supports three apply modes:
\begin{itemize}[nosep]
  \item \code{serial}: a single-threaded Numba level sweep;
  \item \code{parallel}: a \code{prange} loop inside each level;
  \item \code{auto}: select \code{parallel} only when the maximum level width
        is above a configurable threshold.
\end{itemize}
For the current \code{test2} diagnostics, the serial Numba sweep is often the
better baseline because level widths are in the hundreds rather than the
thousands.

\section{Diagnostics}

The builder reports:
\begin{itemize}[nosep]
  \item number of edge blocks;
  \item trace block size;
  \item number of topological levels;
  \item maximum, median, and mean level width;
  \item retained upstream block couplings;
  \item dropped same-level block couplings;
  \item dropped downstream block couplings;
  \item fraction of off-diagonal block couplings dropped;
  \item selected apply mode;
  \item setup timings for CSR preparation, coupling counting, block filling,
        diagonal-block inversion, and warm-up.
\end{itemize}

The dropped-coupling counts are important.  A high same-level count means the
upwind graph levels are not sufficient to represent the matrix-level
dependencies exactly.  In that case the preconditioner is still valid, but it
is more like a level block-Jacobi/Gauss-Seidel hybrid.

\section{Benchmark Entry Points}

The current benchmark script adds two configurations:
\begin{lstlisting}
scipy_bicgstab_upwind_bgs
scipy_gmres_upwind_bgs
scipy_bicgstab_upwind_fbgs
scipy_gmres_upwind_fbgs
\end{lstlisting}

They live in \code{scripts/advection_reaction/benchmark_adv_rea_solvers.py}.  The script assembles
the trace matrix once, builds the reusable CSR matrix once, and then compares
the new preconditioner with SciPy ILU and PETSc ILU/ASM configurations.
When multiple SciPy configurations use the same upwind block-GS options, the
benchmark caches and reuses the scaled matrix and preconditioner.  This avoids
charging the setup cost separately to BICGSTAB and GMRES.

\section{Next Optimization}

If the CSR-built version is competitive, the next optimization is to build the
same block-CSR preconditioner data during the Numba trace assembly.  That
should reduce memory traffic and avoid scanning the scalar CSR matrix.  The
apply kernel should remain in \code{hdgfem.linalg.upwind_block_gs}; only the
builder should move closer to the fused assembly path.

\end{document}
